\section{The compact family and its period coordinates}\label{sec:geometry}

\subsection{The central transition}

Write $B_0=\PP^1_f\times\PP^1_g$, $L=\OO_{B_0}(1,-1)$, and
use the quotient convention in \eqref{eq:ambient-intro}. Let
$\xi_1,\xi_2$ be the two tautological divisor classes and put
$k_i=g+\xi_i$. Each $k_i$ is globally generated, since the direct
image of its line bundle is $\OO(0,1)\oplus\OO(1,0)$.
If $[s_0:s_1],[t_0:t_1]$ are coordinates on $B_0$, the resulting
map to $A=\PP^3\times\PP^3$ has local formula
\begin{equation}\label{eq:central-map}
 (x;y)=(v_1s;u_1t),\qquad
 (x';y')=(v_2s;u_2t),
\end{equation}
where $(u_i,v_i)$ are the quotient coordinates. This formula respects
the changes of trivialization of $L$, and pulls back $\OO_A(1,0)$
and $\OO_A(0,1)$ to $k_1$ and $k_2$.

Its image is the complete intersection
\begin{equation}\label{eq:central-fourfold}
 V_0=\{x_0x'_1-x_1x'_0=0,\quad
         y_0y'_1-y_1y'_0=0\}\subset A.
\end{equation}
Indeed the two determinant equations specify the common directions
$s,t$ whenever the corresponding pairs are nonzero. One direction is
free exactly when both its pairs vanish; both cannot be free at the
same point of $A$. The remaining scalars then reconstruct
\eqref{eq:central-map}. This proves set-theoretic surjectivity and
irreducibility of the image. The two equations form a regular sequence
and are generically smooth, so their complete intersection is reduced
and has the asserted scheme structure.

The singular locus consists of the two disjoint lines
\begin{equation}\label{eq:singular-lines}
 \ell_x=\{x=x'=0,\ [y]=[y']\},\qquad
 \ell_y=\{y=y'=0,\ [x]=[x']\}.
\end{equation}
Near $\ell_x$, use nonzero $y,y'$ coordinates and eliminate the
second equation. The first is a $2\times2$ determinant in the four
$x$ coordinates, with a free coordinate along the line. The same
argument applies to $\ell_y$. Thus the transverse singularities are
ordinary threefold nodes. The complete intersection is Cohen--Macaulay
and regular in codimension one, hence normal.

The exceptional loci of $Z\to V_0$ are the two sections
$S_f,S_g\simeq B_0$ on which both quotients are respectively onto
$\OO$ and onto $L$. Each contracts one ruling to its corresponding
line. Their codimension in $Z$ is two, so the map is small.
Adjunction and the projective-bundle formula give
\begin{equation}\label{eq:central-canonical}
 \omega_{V_0}=\OO_{V_0}(-2,-2),\qquad
 -K_Z=4g+2\xi_1+2\xi_2=2k_1+2k_2.
\end{equation}
The contraction is therefore crepant. Normality gives
$q_{Z*}\OO_Z=\OO_{V_0}$, and projection formula identifies its
anticanonical section ring with that of the ample line bundle
$\OO_{V_0}(2,2)$. In particular $V_0=\bar Z$ is the anticanonical
model in the central construction.

Choose a general anticanonical $Y\subset Z$, also nondegenerate for
the toric polytope, and let $X_0$ be its image. The defining section
of $\OO_{V_0}(2,2)$ lifts to $A$: the twisted two-equation Koszul
resolution is
\[
 0\longrightarrow\OO_A\longrightarrow\OO_A(1,1)^{\oplus2}
      \longrightarrow I_{V_0}(2,2)\longrightarrow0,
\]
and its higher cohomology vanishes. Hence $X_0$ is also described
by the two equations in \eqref{eq:central-fourfold} and one equation
of bidegree $(2,2)$. Its section meets each line in
\eqref{eq:singular-lines} in four simple points.

The same multiplicities and the exceptional periods can be computed
on $Z$. The normal bundles and anticanonical restrictions of the two
sections are
\begin{equation}\label{eq:section-normals}
 N_{S_f/Z}=(L^{-1})^{\oplus2},\quad N_{S_g/Z}=L^{\oplus2},\qquad
 (-K_Z)|_{S_f}=\OO(0,4),\quad
 (-K_Z)|_{S_g}=\OO(4,0).
\end{equation}
The all-$\OO$ and all-$L$ summands of
\[
 \pi_*\OO_Z(-K_Z)=\OO(0,4)\otimes
       \operatorname{Sym}^2(\OO\oplus L)\otimes
       \operatorname{Sym}^2(\OO\oplus L)
\]
have independent coefficients. The restrictions are consequently the
complete quartic systems. At a simple zero the normal derivative of
the equation removes the trivial summand in the normal bundle of the
contracted ruling. Each resulting curve has normal bundle
$\OO(-1)\oplus\OO(-1)$ in $Y$.

The global toric calculation in \cite[Appendix~A]{wuebben-chambers}
gives $h^{1,1}(Y)=4$, $h^{2,1}(Y)=60$ and the spanning divisor
classes $(f,g,\xi_1,\xi_2)$. Simple connectivity follows from the
toric lattice-quotient calculation in \cite[Section~8]{wuebben-hodge}.
We use these established global data. In these coordinates the two
curve types have degrees
\begin{equation}\label{eq:curve-types}
 \Gamma_f=(1,0,0,0),\qquad \Gamma_g=(0,1,-1,-1).
\end{equation}
The contraction described above has no other exceptional locus. The
polarization retained in this paper is
\begin{equation}\label{eq:polarization}
 \cL_0=\OO_{V_0}(2,2)|_{X_0},\qquad
 H=c_1(q^*\cL_0)=(4g+2\xi_1+2\xi_2)|_Y.
\end{equation}
In particular it is not replaced by a multiple when an embedding is
chosen below.

\subsection{Local periods through a node}

\begin{lemma}\label{lem:local-period}
Let $Q(z)=z_1^2+\cdots+z_4^2$ and let $Q(z)=\tau(s)$ be a
holomorphic family of ordinary double points. For a nowhere-zero
relative dualizing form $\Omega$, the oriented local vanishing
period is a single-valued holomorphic function
$P(s)=\tau(s)u(s)$, with $u$ holomorphic and $u(0)\ne0$.
\end{lemma}
\begin{proof}
Write
\[
 \Omega=a(z,s)\operatorname{Res}
       \frac{\dd z_1\wedge\cdots\wedge\dd z_4}{Q(z)-\tau(s)},
 \qquad a(0,0)\ne0,
\]
choosing a holomorphic ambient representative for $a$. For positive
real $\tau$, the real vanishing sphere has radius $\sqrt\tau$.
Scaling $z=\sqrt\tau\,y$ multiplies the residue form by $\tau$.
With the radial residue orientation its integral is $\pi^2\tau$:
the residue is $(2\sqrt\tau)^{-1}$ times the induced sphere volume,
and the sphere volume is $2\pi^2\tau^{3/2}$.

Expand $a$ in its convergent Taylor series in $z$. Odd homogeneous
terms integrate to zero under the antipodal map of $S^3$, which
preserves orientation. A term of degree $2k$ contributes $\tau^{k+1}$
times a holomorphic function of the remaining parameters. Uniform
convergence on a smaller coordinate and parameter domain therefore
gives $P=\tau u$, with $u(0)=\pm\pi^2a(0,0)$. The two choices
of square root give the same integral by the same antipodal symmetry.
Analytic continuation is consequently single-valued, and the formula
extends holomorphically across $\tau=0$.
\end{proof}

\subsection{The six-dimensional polarized smoothing}

\begin{proposition}\label{prop:period-family}
There is a proper flat family $\pi:\mathcal X\to B$ over a ball
$B\subset R$ as in \eqref{eq:period-base-intro}, projective over
$B$, whose central fibre is $X_0$. It has a relatively ample line
bundle restricting to the exact $\cL_0$ and a nowhere-zero relative
holomorphic three-form. Its total space is smooth. The local equations,
oriented periods and bounded units satisfy
\eqref{eq:node-unit-intro}. Its complete discriminant is
$\bigcup_i\{p_i=0\}$; every smooth fibre is simply connected and
has $b_2=2$.
\end{proposition}
\begin{proof}
The central variety is Gorenstein Calabi--Yau, with rational local
complete-intersection singularities. Rationality of the nodes and
the Calabi--Yau resolution give
$H^1(\OO_{X_0})=H^2(\OO_{X_0})=0$. Its singularities are Du~Bois
and the projective resolution satisfies the
$\partial\bar\partial$ lemma. Friedman's unobstructedness theorem
\cite[Theorem~1.6]{friedman-unobstructed} therefore gives a smooth
abstract deformation germ.

Choose an integer $N$ such that $\cL_0^N$ is very ample and
$H^1(\cL_0^N)=0$. Over a small Artin extension,
$H^2(\OO_{X_0})=0$ lifts the embedding line bundle and
$H^1(\cL_0^N)=0$ lifts its sections. The lifted morphism remains
a closed immersion over the nilpotent thickening. Thus embedded
deformations map formally smoothly to abstract deformations. The
restricted Euler sequence expresses the same fact as
$H^1(T_{\PP^M}|_{X_0})=0$. A neighbourhood $S$ of the central
point in the Hilbert germ is smooth and its tangent map onto the
abstract deformation tangent is surjective. Shrink it to retain
integral fibres.

The use of $\cL_0^N$ for this embedding does not change the prescribed
polarization. After an étale neighbourhood, the universal family has
a section through a smooth central point \cite[Tag~055S]{stacks}.
The relative Picard scheme and its Poincaré bundle normalized along
this section are available by
\cite[Theorems~2.5(2), 4.8; Exercises~3.11, 4.3 and their solutions]
{kleiman-picard}. Since $H^2(\OO_{X_0})=0$, the Picard morphism is
smooth at $[\cL_0]$ \cite[Proposition~5.19]{kleiman-picard}.
A further étale neighbourhood has a section through that point.
Pulling back the normalized Poincaré bundle extends $\cL_0$ itself,
and it is relatively ample after shrinking. Both base changes preserve
the tangent image.

The parametric holomorphic Morse lemma gives node functions
$\tau_i:S\to\CC$. The local-to-global tangent calculation of
Rollenske--Thomas \cite[Section~3, (3.2)]{rollenske-thomas} identifies
the image of their joint differential, up to nonzero choices of local
generators, with the kernel of the exceptional curve-class map. By
\eqref{eq:curve-types} this kernel has dimension six and contains a
vector with every coordinate nonzero. Each $\dd\tau_i(0)$ is therefore
nonzero. The total space is smooth at the nodes and elsewhere near
the central fibre. Properness permits a shrink making it smooth
throughout.

The flat family is Gorenstein. Its relative dualizing line bundle
restricts to $\omega_{X_0}\simeq\OO_{X_0}$. Since
$H^1(\omega_{X_0})=0$, cohomology and base change lifts a nonzero
central generator to a relative section. Its zero set has proper
image missing the origin. A further shrink gives a nowhere-zero
relative form $\Omega$.

We next determine the exact, rather than first-order, period relations.
Delete disjoint node balls to obtain a common exterior $U$. Fixed
collars and the proper submersion theorem with boundary identify its
nearby exteriors, including parameters with some singular caps. The
Thom and Mayer--Vietoris sequences give
\begin{equation}\label{eq:attachment-geometry}
 H^2(Y;\ZZ)\xrightarrow{r}\ZZ^8
       \xrightarrow{\partial}H^3(X_s;\ZZ),\qquad
 r(D)_i=D\cdot C_i,\quad
 \partial e_i=\operatorname{PD}_L(\delta_i).
\end{equation}
Here $\operatorname{PD}_L$ uses the oriented Thom class of the
vanishing sphere. To identify the maps, the exceptional curves have
real codimension four, so restriction gives $H^2(Y)\simeq H^2(U)$.
A resolution cap retracts to $S^2$, and restriction to its link records
the curve degree. A smoothing cap retracts to $S^3$ and has no
degree-two cohomology. Its boundary generator maps to the Thom dual
of that zero section. These observations give exactness and the
displayed maps.

The restriction vectors of $f,g$ are the two block indicators
$v_f,v_g$. Integral Poincaré duality therefore gives
\begin{equation}\label{eq:integral-period-relations}
 \sum_{i\in I_f}\delta_i=0,\qquad
 \sum_{i\in I_g}\delta_i=0
       \quad\hbox{in }H_3(X_s;\ZZ).
\end{equation}
The four divisor classes span rational $H^2(Y)$, so these span all
rational relations. No assertion about saturation of the vanishing
image is needed.

Orient the local spheres compatibly with
\eqref{eq:attachment-geometry}. Lemma~\ref{lem:local-period}
gives holomorphic periods $P_i=\tau_i u_i$. On the smooth locus,
\eqref{eq:integral-period-relations} gives
$\sum_{I_f}P_i=\sum_{I_g}P_i=0$. The smooth locus is dense:
the smooth base integrates a full-support smoothing tangent, or,
equivalently, the finite union of the nontrivial hypersurfaces
$\{\tau_i=0\}$ has empty interior. Holomorphic extension gives
these identities on all of $S$.

Consequently $P=(P_i)$ maps $S$ to $R$, and
\[
 \dd P(0)=\operatorname{diag}(u_i(0))\dd\tau(0)
\]
has rank six. The holomorphic submersion theorem gives a section
$\sigma:B\to S$, with $\sigma(0)=0$ and $P\sigma=\mathrm{id}$.
Pull back the compact family, polarization and relative form along
$\sigma$. Lemma~\ref{lem:local-period} now gives
$p_i=\tau_i(p)u_i(p)$. After shrinking the ball, all units are
bounded above and away from zero.

Properness excludes new singularities outside the chosen node charts.
Inside chart $i$, the fibre is singular exactly when $\tau_i=0$,
equivalently $p_i=0$. On that hyperplane
$\dd\tau_i=u_i^{-1}\dd p_i\ne0$, since the coordinate functional
$p_i$ is nonzero on $R$. The pulled-back total space is therefore
smooth as well. The discriminant need not be a normal-crossings divisor
at intersections of its hyperplanes.

Finally the links $S^2\times S^3$ and both kinds of fillings are
simply connected. Van Kampen identifies the fundamental group of a
smooth fibre with that of $Y$, which is trivial. The exceptional
curve map has rank two, so the surgery sequence gives
$b_2(X_p)=b_2(Y)-2=2$. This proves all assertions.
\end{proof}

\subsection{Intersection data and permitted scales}

The Chow relations and normalization are
\[
 f^2=g^2=0,\qquad \xi_i^2=(f-g)\xi_i,\qquad
 \int_Zfg\xi_1\xi_2=1.
\]
Integrating $D_iD_j(-K_Z)^2$ gives
\begin{equation}\label{eq:intersection-matrix}
 B_{ij}:=\int_YD_iD_jH=
 \begin{pmatrix}
 0&8&4&4\\8&0&12&12\\4&12&-8&0\\4&12&0&-8
 \end{pmatrix},\qquad (D_i)=(f,g,\xi_1,\xi_2),
 \qquad H^3=320.
\end{equation}
The kernel of restriction is spanned by
$k_1=g+\xi_1$, $k_2=g+\xi_2$, and $H=2k_1+2k_2$.
Relative to these surviving classes and the quotient lifts $f,g$,
the intersection blocks are
\begin{equation}\label{eq:intersection-blocks}
 B_{KK}=\begin{pmatrix}16&24\\24&16\end{pmatrix},\quad
 B_{DK}=\begin{pmatrix}12&12\\12&12\end{pmatrix},\quad
 B_{DD}=\begin{pmatrix}0&8\\8&0\end{pmatrix}.
\end{equation}
The primitive Hodge pairing on the orthogonal quotient is consequently
\begin{equation}\label{eq:mass-computation}
 -B_{DD}+B_{DK}B_{KK}^{-1}B_{KD}
       =\frac45\begin{pmatrix}9&-1\\-1&9\end{pmatrix}.
\end{equation}
The analytic identification of this pairing with the limiting quotient
mass uses harmonic representatives and removal of all surviving
classes. The intersection calculation itself introduces no chosen
actual small eigenform.

The two-parameter slice $(a,a,a,-3a;b,b,b,-3b)$ lies in $R$.
Every pair of sufficiently small nonzero $a,b$ gives a smooth fibre,
with no comparison condition on their magnitudes. Boundedness of the
units gives
\begin{equation}\label{eq:log-unit-control}
 \bigl|\log|\tau_i(p)|-\log|p_i|\bigr|\le C.
\end{equation}
For $a=e^{-F(h)}$, $b=e^{-G(h)}$, where $F,G\to\infty$, this yields
$w_f\sim8\pi/F$ and $w_g\sim8\pi/G$. It also explains why finite
vanishing orders in one holomorphic arc behave differently: if
$\tau_i(s)=c_i s^{\nu_i}+O(s^{\nu_i+1})$, then
\[
 \kappa_i(s)\sim\frac{2\pi}{\nu_i|\log|s||}.
\]
All such capacities have one logarithmic order. The exact divisibility
$\tau_i=p_i/u_i$ is what permits the more separated paths without
uncontrolled higher-order additive terms in their smallest coordinates.
